%------------------------------------------------------------------------------%

% 01
%------------------------------------------------------------------------------%
\begin{question}
  Calcule o vetor tangente a curva
  \[
    x = t\cos(t)  \qquad
    y = t\sin(t)  \qquad
    0 \leq t < \infty
  \]
  quando \(t=\frac{\pi}{2}\)
\end{question}

%------------------------------------------------------------------------------%
\begin{resolution}{Solução}
  \begin{columns}
    \column{0.5\linewidth}
    \onslide<+->{Derivada da coordenada $x$}
    \begin{align*}
      \onslide<+->{x'(t) & = \frac{dx}{dt}                \\}
      \onslide<+->{      & = \frac{d}{dt}t\cos(t)         \\}
      \onslide<+->{      & = 1\times\cos(t) + t(-\sin(t)) \\}
      \onslide<+->{      & = \cos(t) - t\sin(t)             }
    \end{align*}

    \column{0.5\linewidth}
    \onslide<+->{Derivada da coordenada $y$}
    \begin{align*}
      \onslide<+->{y'(t) & = \frac{dy}{dt}              \\}
      \onslide<+->{      & = \frac{d}{dt}t\sin(t)       \\}
      \onslide<+->{      & = 1\times \sin(t) + t\cos(t) \\}
      \onslide<+->{      & = \sin(t) + t\cos(t)           }
    \end{align*}
  \end{columns}
\end{resolution}

%------------------------------------------------------------------------------%
\begin{resolution}{Solução}
  \onslide<+->{Avaliando as funções derivada em \(t=\frac{\pi}{2}\)}
  \[
    \onslide<+->{x'\left(\frac{\pi}{2}\right)}
    \onslide<+->{   = \left(\cos(t) - t\sin(t)\right) \evalat{\frac{\pi}{2}}{}     }
    \onslide<+->{   = \cos\left(\frac{\pi}{2}\right) - \frac{\pi}{2}\sin\left(\frac{\pi}{2}\right)}
    \onslide<+->{   = 0 - \frac{\pi}{2} \times 1}
    \onslide<+->{   = -\frac{\pi}{2}}
  \]
  \[
    \onslide<+->{y'\left(\frac{\pi}{2}\right)}
    \onslide<+->{   = \left(\sin(t) + t\cos(t)\right) \evalat{\frac{\pi}{2}}{}     }
    \onslide<+->{   = \sin\left(\frac{\pi}{2}\right) + \frac{\pi}{2}\cos\left(\frac{\pi}{2}\right)}
    \onslide<+->{   = 1 + \frac{\pi}{2} \times 0}
    \onslide<+->{   = 1}
  \]

  \vspace{\baselineskip}
  \onslide<+->{O vetor tangente é
    \(
    \begin{bmatrix}
      \sfrac{-\pi}{2} \\ 1
    \end{bmatrix}
    \)}
\end{resolution}



% 02
%------------------------------------------------------------------------------%
\begin{question}
  Calcule o vetor tangente a curva
  \[
    x = \sec(t)  \qquad
    y = \tan(t)  \qquad
    -\frac{\pi}{2} < t < \frac{\pi}{2}
  \]
  no ponto \(t=\frac{\pi}{4}\)
\end{question}

%------------------------------------------------------------------------------%
\begin{resolution}{Solução}
  \begin{columns}
    \column{0.5\linewidth}
    \onslide<+->{Derivada da coordenada $x$}
    \begin{align*}
      \onslide<+->{x'(t) & = \frac{dx}{dt}       \\[2mm]}
      \onslide<+->{      & = \frac{d}{dt}\sec(t) \\[2mm]}
      \onslide<+->{      & = \tan(t)\sec(t)      }
    \end{align*}

    \column{0.5\linewidth}
    \onslide<+->{Derivada da coordenada $y$}
    \begin{align*}
      \onslide<+->{y'(t) & = \frac{dy}{dt}       \\[2mm]}
      \onslide<+->{      & = \frac{d}{dt}\tan(t) \\[2mm]}
      \onslide<+->{      & = \sec^2(t)           }
    \end{align*}
  \end{columns}
\end{resolution}

%------------------------------------------------------------------------------%
\begin{resolution}{Solução}
  \onslide<+->{Avaliando as derivadas em \(t=\frac{\pi}{4}\)}
  \[
    \onslide<+->{x'\left(\frac{\pi}{4}\right)}
    \onslide<+->{   = \tan\left(\frac{\pi}{4}\right)\sec\left(\frac{\pi}{4}\right)}
    \onslide<+->{   = 1 \times \sqrt{2}}
    \onslide<+->{   = \sqrt{2}}
  \]
  \[
    \onslide<+->{y'\left(\frac{\pi}{4}\right)}
    \onslide<+->{   = \sec^2\left(\frac{\pi}{4}\right)}
    \onslide<+->{   = \left(\sqrt{2}\right)^2}
    \onslide<+->{   = 2}
  \]

  \vspace{\baselineskip}
  \onslide<+->{O vetor tangente é
    \(
    \begin{bmatrix}
      \sqrt{2} \\ 2
    \end{bmatrix}
    \)}
\end{resolution}



% 03
%------------------------------------------------------------------------------%
\begin{question}{Exemplo \theexemplo}
  Calcule o vetor tangente a curva
  \[
  \begin{cases}
    x & = e^{2t} \cos(\pi t) \\
    y & = e^t \sin(\pi t)
  \end{cases}
  \]
  em \(t=\frac{3}{4}\)
\end{question}

%------------------------------------------------------------------------------%
\begin{resolution}{Solução}
  \onslide<+->{Derivada da coordenada $x$}
  \begin{align*}
    \onslide<+->{x'(t) & = \frac{dx}{dt}       \\}
    \onslide<+->{      & = \frac{d}{dt}\left(e^{2t} \cos(\pi t)\right) } \\
    \onslide<+->{      & = \frac{d}{dt}\left(e^{2t}\right) \cos(\pi t)
      + e^{2t} \frac{d}{dt}\left(\cos(\pi t)\right) } \\
    \onslide<+->{      & = e^{2t}\frac{d}{dt}\left(2t\right) \cos(\pi t)
      - e^{2t} \sin(\pi t)\frac{d}{dt}\left(\pi t\right) } \\
    \onslide<+->{      & = 2 e^{2t} \cos(\pi t) - \pi e^{2t} \sin(\pi t) }
  \end{align*}
\end{resolution}

%------------------------------------------------------------------------------%
\begin{resolution}{Solução}
  \onslide<+->{Derivada da coordenada $y$}
  \begin{align*}
    \onslide<+->{y'(t) & = \frac{dy}{dt}       \\[2mm]}
    \onslide<+->{      & = \frac{d}{dt}\left(e^t \sin(\pi t)\right)} \\
    \onslide<+->{      & = \frac{d}{dt}\left(e^t\right) \sin(\pi t)
      + e^t \frac{d}{dt}\left(\sin(\pi t)\right)} \\
    \onslide<+->{      & = e^t \sin(\pi t)
      + e^t \cos(\pi t)\frac{d}{dt}\left(\pi t\right)} \\
    \onslide<+->{      & = e^t \sin(\pi t)+ \pi e^t \cos(\pi t)}
  \end{align*}
\end{resolution}

%------------------------------------------------------------------------------%
\begin{resolution}{Solução}
  \onslide<+->{Avaliando a derivada de $x(t)$ em \(t=\sfrac{3}{4}\)}
  \begin{align*}
    \onslide<+->{x'\left(\frac{3}{4}\right)}
    \onslide<+->{& = \left(2 e^{2t} \cos(\pi t) - \pi e^{2t} \sin(\pi t)\right)
      \evalat{t=\sfrac{3}{4}}{}}
    \\
    \onslide<+->{& =   2 e^{2\frac{3}{4}} \cos\left(\pi \frac{3}{4}\right)
      - \pi e^{2\frac{3}{4}} \sin\left(\pi \frac{3}{4}\right)}
    \\
    \onslide<+->{& =   2 e^{\sfrac{3}{2}} \cos\left(\frac{3\pi}{4}\right)
      - \pi e^{\sfrac{3}{2}} \sin\left(\frac{3\pi}{4}\right)}
    \\
    \onslide<+->{& =   2 e^{\sfrac{3}{2}} \left(\frac{-\sqrt{2}}{2}\right)
      - \pi e^{\sfrac{3}{2}} \frac{\sqrt{2}}{2}}
    \\
    \onslide<+->{& = - \sqrt{2} \left(1 + \frac{\pi}{2} \right) e^{\sfrac{3}{2}}}
  \end{align*}
\end{resolution}

%------------------------------------------------------------------------------%
\begin{resolution}{Solução}
  \onslide<+->{Avaliando a derivada de $y(t)$ em \(t=\sfrac{3}{4}\)}
  \begin{align*}
    \onslide<+->{y'\left(\frac{3}{4}\right)}
    \onslide<+->{& = \big(e^t \sin(\pi t)+ \pi e^t \cos(\pi t)\big)
      \evalat{t=\sfrac{3}{4}}{}}
    \\
    \onslide<+->{& =     e^{\sfrac{3}{4}} \sin\left(\frac{3\pi}{4}\right)
      + \pi e^{\sfrac{3}{4}} \cos\left(\frac{3\pi}{4}\right)}
    \\
    \onslide<+->{& =     e^{\sfrac{3}{4}} \frac{\sqrt{2}}{2}
      - \pi e^{\sfrac{3}{4}} \frac{\sqrt{2}}{2}}
    \\
    \onslide<+->{& = \sqrt{2} \frac{1 - \pi}{2} e^{\sfrac{3}{4}}}
  \end{align*}
\end{resolution}

%------------------------------------------------------------------------------%
\begin{resolution}{Solução}
  O vetor tangente no ponto \(t=\sfrac{3}{4}\) é
  \[
    \frac{d}{dt}\left(\!
    \begin{array}{c}
      x \\ y
    \end{array}
    \!\right)
    =
    \left(
    \begin{array}{c}
      - \sqrt{2} \left(1 + \dfrac{\pi}{2} \right) e^{\sfrac{3}{2}} \\[9mm]
      \sqrt{2} \; \dfrac{1 - \pi}{2} e^{\sfrac{3}{4}}
    \end{array}
    \right)
  \]
\end{resolution}

%------------------------------------------------------------------------------%
